\subsection{An invertible adjacency operator on all functions: Problem 15.89}
\label{cand:15.89}
\problemstatement{15.89}

Let \(G=C_3*C_2=\langle a,b:a^3=b^2=1\rangle\) and
\[
 S=\{aba,aba^2,a^2ba,a^2ba^2\}.
\]
Its Cayley graph, with edges \(g\sim gs\), is simple
and four-regular: free-product normal form gives four
distinct nonidentity words, and \(S=S^{-1}\).
It is connected because, for \(s_1=aba,s_2=a^2ba\),
\(s_1s_2^{-1}=a^2\) and \(a^2s_1a^2=b\).
The group is infinite by the distinct reduced words
\((ab)^n\), and left translation is vertex-transitive.

On the space \(\C^G\) of all functions, put
\((R_gf)(x)=f(xg)\), \(U=R_a+R_{a^2}\), and
\(B=R_b\). Then \(R_gR_h=R_{gh}\), and
\[
 U^2=U+2I,\qquad U^{-1}=(U-I)/2,\qquad B^2=I.
\]
The adjacency operator is exactly \(A=UBU\), hence
\[
 A^{-1}=\frac14(U-I)B(U-I).
\]
Both products with \(A\) equal the identity.
All these expressions are finite sums of translations,
so they act on arbitrary, even unbounded, functions.
There is no support or convergence condition.
In particular zero is not an eigenvalue, answering
the question negatively.

The factorization has the familiar zig-zag form:
a move within a triangle, across a matching edge,
and within the next triangle. The general operation
is due to \cite{zigzag}; the invertibility of finite
factors also occurs in \cite{permutational-powers}.
The proposed contribution is this particular infinite
simple vertex-transitive example and its use for
the printed question, not the graph-product operation.
Source: \artifact{research/15.89-proof.md}.

